\documentclass[11pt,a4paper]{article}
\usepackage[utf8]{inputenc}
\usepackage[T1]{fontenc}
\usepackage{lmodern}
\usepackage[margin=0.95in]{geometry}
\usepackage{amsmath,amssymb,amsthm,mathtools}
\usepackage{booktabs,graphicx,array,microtype,xcolor}
\usepackage{hyperref}
\hypersetup{colorlinks=true,linkcolor=blue!60!black,citecolor=blue!60!black,urlcolor=blue!60!black}
\setlength{\emergencystretch}{3em}
\makeatletter
\DeclareRobustCommand{\_}{\ifmmode\nfss@text{\textunderscore}\else\textunderscore\allowbreak\fi}
\makeatother
\newtheorem{theorem}{Theorem}[section]
\newtheorem{proposition}[theorem]{Proposition}
\newtheorem{lemma}[theorem]{Lemma}
\newtheorem{corollary}[theorem]{Corollary}
\theoremstyle{definition}
\newtheorem{definition}[theorem]{Definition}
\newtheorem{remark}[theorem]{Remark}

\title{Parallel balanced minimum evolution optimization\\in Phylo2Vec coordinates}
% Author names and affiliations to be supplied.
\author{}
\date{}

\begin{document}
\maketitle

\begin{abstract}
Balanced minimum evolution (BME) provides a distance-based criterion for phylogenetic inference, but searching large tree spaces requires repeated evaluation of many candidate rearrangements.
We formulate these evaluations as discrete coordinate differences in Phylo2Vec space and organize the underlying subtree-distance recurrences for parallel execution.
For a tree represented in an ordered labeling, all valid single-coordinate changes are evaluated in $O(N^2)$ time.
Relabeling preserves the current tree while changing the coordinate neighborhood.
To reduce memory requirements, we optimize smaller skeleton trees whose leaves represent fixed components of the full tree.
We give a self-contained proof that dyadically weighted component distances preserve every skeleton-induced BME change exactly.
This reduction also covers a subtree window when its entire complement is retained as an oriented component.
Capped profiles accelerate proposal generation, with full-profile evaluation used for acceptance and outside profiles refreshed after each committed window.
In 120 small-tree comparisons, the CPU optimizer matches or improves the reference BME score in 119 cases.
A full gradient sweep at $N=40{,}000$ takes 158.36~ms on eight H100 GPUs.
On a 100,000-sequence RNASim benchmark, minibatch refinement with full-profile acceptance takes 104.45~s and returns 8,802 incorrect splits, compared with 593.50~s and 8,857 splits for VeryFastTree refinement.
Exact evaluation applies to the specified distance matrix and candidate moves, while the search remains heuristic.
\end{abstract}

\section{Introduction}
\label{sec:introduction}

Distance-based phylogenetic inference separates the estimation of pairwise evolutionary distances from the search for a tree.
Balanced minimum evolution (BME) scores a topology using a weighted sum of those distances, with weights determined by the branching structure along each leaf-to-leaf path~\cite{pauplin2000direct,desper2002fast,desper2004theory}.
Efficient recurrences permit an entire subtree-prune-and-regraft (SPR) neighborhood to be evaluated in quadratic time~\cite{hordijk2005spr,lefort2015fastme}.
For large trees, the practical costs include storing the subtree-distance table, repeatedly traversing the tree, and coordinating many small candidate evaluations.
The number of optimization sweeps depends on the instance and the search strategy.

Phylo2Vec represents rooted binary trees by bounded integer vectors~\cite{penn2024phylo2vec}.
This representation provides a coordinate system in which a discrete difference records the change in BME after replacing one vector entry.
An ordered labeling permits efficient decoding and maps coordinate changes to a structured collection of SPR moves.
The representation also separates the tree topology from its leaf-label permutation.
Changing the permutation can expose different coordinate moves without changing the current labeled tree or its score.

Our method combines this coordinate formulation with parallel evaluation and component-based reduction.
The reduction relies on established BME structure.
Balanced subtree averages underlie earlier distance-based algorithms~\cite{desper2002fast,desper2004theory}, and fixed-clade families have been described as affine images of smaller BME polytopes~\cite[Section~2.4]{haws2011optimality}.
We state the corresponding objective identity explicitly because its assumptions determine when minibatch optimization evaluates the full-tree objective.
In particular, a component must retain its complete balanced leaf profile, and a subtree window must retain the rest of the tree through the correct attachment orientation.

The contribution developed here is an algorithmic formulation: a quadratic-time Phylo2Vec coordinate evaluator, an optimizer that verifies proposed updates against BME, and parallel implementations using column decomposition and reduced skeletons.
We distinguish exact identities from finite-precision calculations and from heuristic search decisions.
The experiments separately assess objective agreement, local optimization, gradient throughput, and sequence-tree refinement.

\section{Methods}
\label{sec:math}

\subsection{BME and Phylo2Vec coordinates}

Let $T$ be an unrooted binary tree on $N\geq3$ labeled leaves, and let $D$ be a fixed, finite, symmetric matrix with zero diagonal.
Write $q_T(x,y)$ for the number of internal degree-three vertices on the path from $x$ to $y$.
The BME objective is
\begin{equation}
 L(T;D)=\sum_{0\leq x<y<N}D_{xy}\,2^{-q_T(x,y)}.
 \label{eq:pauplin}
\end{equation}
If the path has $e_T(x,y)$ edges, then $q_T(x,y)=e_T(x,y)-1$.
Thus the unordered-pair coefficient is $2^{1-e_T(x,y)}$.
An artificial degree-two root used to encode the tree is suppressed before evaluating these path lengths.
No additivity assumption on $D$ is needed for the identities below.

The Phylo2Vec domain for rooted binary trees is
\begin{equation}
 \mathcal V_N=\{(v_0,\ldots,v_{N-2}):0\leq v_i\leq2i\},
 \qquad |\mathcal V_N|=(2N-3)!!.
 \label{eq:phylo2vec_space}
\end{equation}
We use the ordered, or cherry-first, subset
$\mathcal V_N^{\mathrm{cf}}=\{v:0\leq v_i\leq i\}$ as the baseline for coordinate evaluation.
Every rooted binary tree can be represented in this subset after a suitable leaf permutation.
The relabeling construction and its proof are given in Appendix~\ref{sec:gauge}.
If $\pi(a)$ is the original leaf represented by current index $a$, the distance matrix in these coordinates is
$D_\pi[a,b]=D[\pi(a),\pi(b)]$.
Permuting both axes in this way preserves the objective.
Constructing the permutation takes $O(N)$ time from a decoded tree, while explicitly copying a dense permuted matrix takes $O(N^2)$ time.

For an ordered baseline $v$, define the discrete coordinate difference
\begin{equation}
 \Delta_{i,j}
 =L\bigl(T(v^{(i\leftarrow j)});D_\pi\bigr)
  -L\bigl(T(v);D_\pi\bigr),\qquad 0\leq j\leq2i.
 \label{eq:discrete_gradient}
\end{equation}
There are $(N-1)^2$ valid coordinate entries, including unchanged coordinates and entries that can represent the same unrooted topology.
This count is not a count of distinct SPR neighbors.
General candidate values $j>i$ are allowed even though the baseline is ordered.
After accepting such a move, the resulting tree is decoded and relabeled before the next ordered evaluation.

\subsection{Quadratic-time coordinate evaluation}

For two distinct vertices $u,w$, define $A(u,w)$ as the balanced average distance between the two components oriented away from the path joining $u$ and $w$.
At an internal vertex, the two outward branches each receive weight $1/2$.
At a leaf, the corresponding distribution is concentrated on that leaf.
Consequently, when the component at $u$ has children $a,b$,
\begin{equation}
 A(u,w)=\tfrac12\{A(a,w)+A(b,w)\}.
 \label{eq:bottom_up}
\end{equation}
The direction is important when one vertex is an ancestor of the other in a computational rooting.
In that case, the component away from $w$ combines the sibling branch and the parent branch:
\begin{equation}
 A(u,w)=\tfrac12\{A(\operatorname{sib}_u(w),w)
                         +A(\operatorname{parent}(u),w)\}.
 \label{eq:top_down}
\end{equation}
At the degree-three computational root, the analogous average uses the two branches not containing $w$.
Bottom-up and top-down passes construct the required table in $O(N^2)$ arithmetic operations.

For a fixed pruned component, the standard BME SPR recurrence propagates score differences across regraft edges in $O(N)$ time~\cite{hordijk2005spr}.
Each propagation step is an adjacent interchange of four directed components.
Its score change is a linear combination of their balanced distances, with coefficient $1/4$.
Successive steps accumulate the change relative to the original attachment.
Phylo2Vec supplies the mapping from each valid coordinate value to one of these regraft edges.
The explicit mapping, including the collapsed-root boundary, is in Appendix~\ref{sec:mapping}.

\begin{proposition}[Cost and scope of a coordinate sweep]
\label{prop:quadratic}
Given an ordered baseline and its leaf permutation, all entries in
Equation~\eqref{eq:discrete_gradient} can be computed in $O(N^2)$ time.
In exact arithmetic, these entries equal the corresponding full-tree BME differences.
\end{proposition}
\begin{proof}
The balanced-distance table costs $O(N^2)$.
For each of $O(N)$ coordinates, pruning its associated component and traversing the two remaining half-trees visits $O(N)$ edges.
The SPR recurrence evaluates the score at each edge, and the coordinate mapping selects the appropriate values without evaluating the candidate trees independently.
The mapping costs at most $O(N)$ per row.
The recurrence preserves the Pauplin score difference at each traversal step, so selecting a mapped edge gives Equation~\eqref{eq:discrete_gradient}.
\end{proof}

The proposition concerns the coordinate neighborhood of the current labeling.
It does not assert that this neighborhood is the entire SPR neighborhood, or that coordinate descent finds a global BME optimum.

\subsection{Exact reduction to a skeleton tree}
\label{sec:condensation}

Partition all $N$ leaves into $K\geq3$ fixed components $S_1,\ldots,S_K$.
Each component has a fixed internal binary topology and a single attachment to the rest of the tree.
A component can consist of one leaf.
For a non-singleton component, let $r_a$ be the internal vertex adjacent to its attachment edge, and let $h_a(x)$ count the branching vertices from $r_a$ to leaf $x$, including $r_a$.
For a singleton, set $h_a(x)=0$.
Define its balanced profile by
\begin{equation}
 \omega_a(x)=2^{-h_a(x)},\qquad x\in S_a.
 \label{eq:profile}
\end{equation}

\begin{lemma}[Profile normalization]
\label{lem:normalization}
For each component, $\sum_{x\in S_a}\omega_a(x)=1$.
\end{lemma}
\begin{proof}
A singleton has weight one.
At a binary root, every leaf weight in each child component is multiplied by $1/2$.
Induction over the component therefore gives total mass
$\tfrac12\cdot1+\tfrac12\cdot1=1$.
\end{proof}

Let $G$ be an unrooted binary skeleton with one leaf for each component.
Construct $T_G$ by replacing skeleton leaf $a$ with $S_a$, retaining the attachment edge.
Define a condensed distance matrix by
\begin{equation}
 \widetilde D_{ab}
 =\sum_{x\in S_a}\sum_{y\in S_b}
    \omega_a(x)\omega_b(y)D_{xy}\quad(a\ne b),
 \qquad \widetilde D_{aa}=0.
 \label{eq:condensed_D}
\end{equation}
This is an objective-level form of the fixed-clade reduction discussed by Haws et al.~\cite[Section~2.4]{haws2011optimality}.
We include the proof to specify the weights and the scope of exactness needed by the algorithm.

\begin{theorem}[BME decomposition over fixed components]
\label{thm:clade_condense}
There is a constant $C$, independent of the skeleton topology, such that
\begin{equation}
 L(T_G;D)=C+L(G;\widetilde D).
 \label{eq:bme_condense_decomp}
\end{equation}
Consequently, for any two skeletons $G$ and $G'$,
\begin{equation}
 L(T_{G'};D)-L(T_G;D)
 =L(G';\widetilde D)-L(G;\widetilde D).
 \label{eq:condensed_delta}
\end{equation}
\end{theorem}
\begin{proof}
Separate the unordered leaf pairs into pairs within a component and pairs between components.
For $x,y\in S_a$, the path stays inside $S_a$ and is unchanged by rearranging the skeleton.
The sum of all such terms is
\[
 C=\sum_a\sum_{\substack{x<y\\x,y\in S_a}}
      D_{xy}\,2^{-q_{T_G}(x,y)},
\]
which is independent of $G$.

For $x\in S_a$ and $y\in S_b$ with $a\ne b$, the path contains the branching vertices inside $S_a$, those on the skeleton path, and those inside $S_b$.
These sets of vertices are disjoint, giving
\[
 q_{T_G}(x,y)=h_a(x)+q_G(a,b)+h_b(y).
\]
Therefore its weight factors as
$2^{-q_{T_G}(x,y)}=\omega_a(x)\,2^{-q_G(a,b)}\,\omega_b(y)$.
Summing between-component terms gives
\begin{align*}
 \sum_{a<b}\sum_{\substack{x\in S_a\\y\in S_b}}
 D_{xy}\,2^{-q_{T_G}(x,y)}
 &=\sum_{a<b}2^{-q_G(a,b)}
      \sum_{\substack{x\in S_a\\y\in S_b}}\omega_a(x)\omega_b(y)D_{xy}\\
 &=L(G;\widetilde D).
\end{align*}
Adding the within-component constant proves Equation~\eqref{eq:bme_condense_decomp}.
Subtracting two such decompositions cancels $C$ and proves Equation~\eqref{eq:condensed_delta}.
\end{proof}

Thus a skeleton evaluates exact changes in the full objective, even though its absolute score omits $C$.
The theorem applies to every rearrangement preserving the chosen components, not only to SPR moves.
It also applies to arbitrary fixed dissimilarities, including non-additive estimated distances.

\begin{corollary}[A proper subtree window]
\label{cor:outside}
For a window attached to the rest of a binary tree by one edge, treat its entire complement as an additional component oriented away from that edge.
If this component retains its full profile and internal topology, Equation~\eqref{eq:condensed_delta} applies to rearrangements of the window skeleton.
\end{corollary}
\begin{proof}
Cutting the attachment edge produces two components.
The endpoint in the complement has two outward branches unless the complement is a singleton.
It therefore defines precisely a binary component of the type used in the theorem.
Including it completes the partition of all leaves, so both within-window and window-to-complement contributions are included.
\end{proof}

The outside profile is determined by the current topology of the complement.
It cannot in general be replaced by one representative leaf, by uniform weights, or by a profile left unchanged after another window has been rearranged.

\subsection{Approximate proposals with full-profile acceptance}
\label{sec:acceptance}

Constructing $\widetilde D$ can be expensive when components contain many leaves.
The sequence implementation therefore permits a cap on the number of profile entries used to propose skeleton moves.
Unexpanded profile mass is assigned to representative leaves.
This defines a surrogate matrix $\widetilde D^{\mathrm{cap}}$.
Without a bound on distances within the omitted subcomponents, a fixed cap gives no uniform accuracy guarantee.
In particular, a balanced component can retain substantial unresolved mass at every frontier entry.

The surrogate selects a candidate skeleton $G'$.
Acceptance then evaluates the full-profile distances for pairs affected by both the topology change and profile truncation.
Uncapped pair distances can be reused.
If $\mathcal P(G,G')=\{(a,b):a<b,\ q_G(a,b)\ne q_{G'}(a,b)\}$, the full change is
\begin{equation}
 \delta=\sum_{(a,b)\in\mathcal P(G,G')}
   \left(2^{-q_{G'}(a,b)}-2^{-q_G(a,b)}\right)
   \widetilde D_{ab}.
 \label{eq:changed_pairs}
\end{equation}
Unchanged pairs contribute zero, and the within-component constant cancels.
Theorem~\ref{thm:clade_condense} therefore identifies $\delta$ with the full-tree BME change.
The cap can affect which proposals are found, but it is not used to determine the final acceptance score.

Windows are committed sequentially.
After a commit, affected ancestor summaries are updated and the next window's outside profile is reconstructed from the current tree.
Node-disjoint windows need not have additive BME changes because both updates can alter the coefficients of distances between their leaves.
For consecutive accepted trees $T^{(0)},\ldots,T^{(M)}$, exact arithmetic gives
\begin{equation}
 L(T^{(M)};D)-L(T^{(0)};D)
 =\sum_{m=1}^M\bigl[L(T^{(m)};D)-L(T^{(m-1)};D)\bigr].
 \label{eq:telescoping}
\end{equation}
If each increment is negative, so is the cumulative increment.
This argument requires that each acceptance calculation uses the current tree, rather than a shared pre-update tree.

In the implementation, changed-pair accumulation uses FP64 and the sequence minibatch path requires a computed improvement greater than $10^{-6}$.
This is a numerical acceptance margin, not a proved bound on floating-point error.
The mathematical equalities are exact in real arithmetic.
Computed distances, profile weights, and reductions remain subject to rounding and, on sufficiently deep trees, underflow.
Exactness here is relative to the specified distance matrix, not to unknown evolutionary distances.

\subsection{Search strategy}
\label{sec:optimizer}

The full-tree CPU optimizer evaluates coordinate differences, ranks improving rows, and forms candidate batches using disjoint path corridors as a conflict filter.
It then evaluates the combined candidate tree and backtracks to a smaller batch if the score does not improve.
The acceptance decision uses the joint objective rather than the sum of individual predicted improvements.
When the current coordinate neighborhood is exhausted, the optimizer changes the ordered labeling and searches again.
Its full-tree configuration also scans remaining unrooted SPR moves before declaring SPR-local convergence.
An optional restart stage applies random perturbations and resumes descent while retaining the best tree found.
Only accepted descent updates and the retained best score are monotone.
The perturbed working trajectory need not be.

The minibatch implementation alternates skeletons at different scales.
A global skeleton rearranges the backbone, local windows expose detailed structure, and overlapping or paired windows propose changes between regions.
Each window follows four steps:
\begin{enumerate}
 \item Select fixed components from the current tree, including an oriented outside component for a proper subtree.
 \item Construct their profiles and a condensed proposal matrix, applying a cap if configured.
 \item Optimize the skeleton and evaluate its proposed change using Equation~\eqref{eq:changed_pairs}.
 \item Commit an improving candidate, graft the components back, and update summaries before selecting or evaluating the next window.
\end{enumerate}
These steps preserve the objective being optimized under the assumptions above.
They do not prove that a particular finite window schedule explores every tree or reaches every improving move.
Nor does an infinite number of batches alone establish such a result without assumptions on which proposals can occur and be accepted.

\subsection{Parallel evaluation and computational cost}
\label{sec:parallel}

Equation~\eqref{eq:bottom_up} has an independent recurrence for each fixed column.
Assigning columns to threads lets neighboring threads access consecutive entries while following the same traversal order.
The top-down correction for column $w$ follows its ancestor path.
Depth-first entry intervals determine the next child without a traversal stack.
Its total work over columns is $O(\sum_w\operatorname{depth}(w))$, bounded by $O(N^2)$ and reduced to $O(N\log N)$ for trees of logarithmic height.

A full balanced-distance table stores approximately $4N^2$ entries.
The column-tiled variant retains the internal-to-leaf block and reconstructs required internal rows on demand.
This reduces that table to approximately $N^2$ entries, while the input matrix and other buffers still contribute to total memory.
A row-wise reduction can return the best coordinate per row without storing the rectangular gradient output.
Columns and SPR rows can also be distributed across GPUs.
Appendix~\ref{sec:parallel_details} gives the dependency and workspace arguments.

For a skeleton with $K$ leaves, a coordinate sweep costs $O(K^2)$ once its distances are available.
This does not make arbitrary distance construction independent of $N$.
Direct evaluation of one condensed pair can require $|S_a||S_b|$ original distances.
Full profiles use $O(N)$ entries, and uncapped changed-pair refinement can dominate runtime for large component pairs.
When $D_{xy}=\phi(x)^\top\psi(y)$ has a fixed-dimensional bilinear representation, weighted component embeddings give
$\widetilde D_{ab}=(\sum_x\omega_a(x)\phi(x))^\top
(\sum_y\omega_b(y)\psi(y))$.
That shortcut is exact for the specified bilinear distance.
It does not justify averaging embeddings before applying a nonlinear sequence-distance correction.

\section{Experimental design}
\label{sec:experiments}

We distinguish four workloads.
First, small-tree tests compare coordinate differences and condensed differences against independent evaluations of Equation~\eqref{eq:pauplin}.
Second, CPU optimization experiments compare final scores and wall times from identical starting trees.
Third, gradient benchmarks measure parallel sweep throughput, independently of convergence.
Fourth, sequence benchmarks measure tree refinement from a shared starting topology.
The large embedding-based scaling workload is distinct from sequence-distance refinement.

The CPU comparison uses additive distances from trees whose branch lengths are uniform on $(0.05,0.55)$ and a noisy variant
$D'_{xy}=D_{xy}\exp(0.15Z_{xy})$, with symmetric standard-normal $Z_{xy}$.
There are respectively 20, 20, 10, 5, and 5 replicates at
$N=20,50,100,150,200$ in each regime.
The reference is the project's in-memory C implementation of the FastME BME-SPR algorithm.
Agreement with the external FastME executable is assessed separately.
A score tolerance of $10^{-9}$ defines wins, ties, and losses.
Single-core timings are reported for Apple M5 and AMD Threadripper PRO 5995WX processors.

GPU evaluation uses NVIDIA RTX A6000 and H100 devices, with additional CPU/Metal measurements in the supplement.
The GPU gradient path uses FP32 working arrays in the reported comparisons, whereas independent score calculations and full-profile acceptance accumulation use FP64.
The sequence comparisons use 64-thread FastME 2.1.6 and VeryFastTree 4.0.5~\cite{lefort2015fastme,pineiro2020veryfasttree}.
VeryFastTree's minimum-evolution refinement is run without its likelihood-refinement phase.
We report refinement time separately from starting-tree construction.
Wrong-split counts measure inferred nontrivial splits absent from the reference tree.
For binary trees on the same leaves, normalized RF distance is this count divided by $N-3$.

Sequence data comprise controlled JC69 simulations and the RNASim benchmark associated with PASTA~\cite{mirarab2015pasta}.
RNASim is simulated sequence data with a known reference tree.
The reported alignments have 1,596 columns.
The minibatch comparison uses gap-aware JC69 distances.
Selected full-tree experiments also use the eigenvalue-based distance approach~\cite{penn2025eigen}, which changes the input objective and is identified separately.
Minibatch timings include full-profile acceptance and sequential reconstruction of live outside profiles.

\section{Results}
\label{sec:results}

\subsection{Agreement with independent objective calculations}

Table~\ref{tab:validation} summarizes the reported numerical checks.
The largest discrepancy in the CPU coordinate tests is $3.553\times10^{-14}$, and the condensation tests agree to $1.564\times10^{-13}$.
These are observed absolute errors on the tested cases, not universal error bounds.
Rooted-vector round trips recover all 1,000 tested vectors at $N=40$.
Unrooted round trips preserve the split sets in a further 1,000 cases.
Exhaustive rooted-tree enumeration at $N=4,\ldots,7$ checks 11,460 encodings.

\begin{table}[tb]
\centering
\caption{Agreement with independent BME calculations. Errors are maximum absolute discrepancies over the stated tests.}
\label{tab:validation}
\small
\begin{tabular}{@{}p{0.39\linewidth}p{0.32\linewidth}r@{}}
\toprule
Calculation & Test scope & Maximum error\\
\midrule
CPU coordinate differences & 36,000 changes, $N=20$ & $3.553\times10^{-14}$\\
Ordered-tree enumeration & 28,488 changes, $N=4$--7 & $2.776\times10^{-15}$\\
Relabeling, first orientation & 200 trees, $N=25$ & $1.776\times10^{-14}$\\
Relabeling, reversed orientation & Same 200 trees & $2.665\times10^{-14}$\\
Condensed SPR differences & 16,231 changes, $N=40$, $K=16$ & $1.564\times10^{-13}$\\
FP32 CUDA gradient & Reported sizes $N=100$--4,000 & $1.11\times10^{-6}$\\
\bottomrule
\end{tabular}
\end{table}

Independent enumeration also identifies the scope of the coordinate neighborhood.
On 30 random trees at each tested size from $N=8$ to $20$, a single ordered orientation covers approximately 27--33\% of distinct SPR neighbors.
Two orientations together cover approximately 52--56\%.
Neither constitutes an exhaustive SPR search.
The full CPU optimizer's additional SPR scan is therefore material: independent enumeration found no improving SPR move in 90 returned trees at $N=10,16,24$.
This observation concerns that tested configuration, not the minibatch variant.

\subsection{CPU optimization}

Both methods recover the reference topology in all 60 noiseless trials.
On the 60 noisy trials, Phylo2Vec obtains two lower BME scores, 57 ties, and one higher score than the in-memory reference.
Thus it matches or improves the reference objective in 119 of 120 cases overall.
A better BME score need not imply a smaller RF distance when the input distances are noisy.
For example, the $N=200$ noisy comparison has a lower mean reference RF error despite one BME-score win for Phylo2Vec.

At $N=200$, the noiseless Apple M5 comparison requires 63.6 gradient evaluations on average, versus 211.0 reference SPR iterations.
The corresponding mean times are 18.14 and 75.46~ms.
Table~\ref{tab:cpu_summary} gives the largest-size timings for both processors.
Supplementary Tables~\ref{tab:bench_noiseless} and~\ref{tab:bench_noisy} give results at all sizes.
External FastME comparisons at the same five sizes produce identical final unrooted splits and BME differences below $10^{-8}$ in the reported checks.
These checks support the reference implementation on those inputs without making the in-memory timings measurements of the external executable.

\begin{table}[tb]
\centering
\caption{Single-core optimization at $N=200$. Times are mean $\pm$ standard deviation in milliseconds over five replicates.}
\label{tab:cpu_summary}
\small
\begin{tabular}{@{}llrrr@{}}
\toprule
Processor & Distances & Reference & Phylo2Vec & Ratio\\
\midrule
Apple M5 & Additive & $75.46\pm2.42$ & $18.14\pm0.74$ & $4.16$\\
Apple M5 & Noisy & $85.16\pm2.06$ & $24.67\pm2.79$ & $3.45$\\
AMD 5995WX & Additive & $160.64\pm5.51$ & $27.96\pm1.18$ & $5.74$\\
AMD 5995WX & Noisy & $181.97\pm4.30$ & $36.91\pm3.94$ & $4.93$\\
\bottomrule
\end{tabular}
\end{table}

\subsection{Gradient throughput and memory}

At $N=4{,}000$, a reported full-gradient evaluation takes 6.219~ms on one H100, compared with 701.375~ms for its CPU reference.
At $N=40{,}000$, distributing the full-table computation across eight H100s reduces time from 816.34 to 158.36~ms, a factor of 5.15.
The SPR stage alone falls from 756.26 to 98.28~ms.
Column tiling reduces storage at the cost of recomputation: the reported one-GPU memory requirement falls from 31,142 to 13,614~MB, while the eight-GPU tiled evaluation takes 1,032.83~ms at this size.
Thus the reduced-memory variant is useful when capacity, rather than latency, is limiting.

Distributed and streamed configurations extend a single full-tree sweep to $N=380{,}000$, taking 205.450~s with 72,946~MB reported per GPU on eight H100s.
This is a sweep benchmark, not the time to infer a converged tree.
For the embedding-based workload with $N=10{,}000{,}000$ and $K=20{,}000$, reported average condensation and gradient times are 10.20 and 48.57~ms per batch on eight H100s.
Multiplying their sum by 667 batches gives the reported epoch estimate of 39.20~s.
It excludes any additional search work not represented by those two timings and does not establish a sequence-refinement time at that scale.
The complete scaling tables are in Appendix~\ref{sec:scaling_tables}.

\subsection{Sequence-tree refinement}

Table~\ref{tab:rnasim_summary} compares minibatch refinement with VeryFastTree refinement from the same initial topology.
At $N=50{,}000$, the minibatch method returns 4,950 incorrect splits in 16.84~s, compared with 4,980 in 313.53~s.
At $N=100{,}000$, the corresponding results are 8,802 in 104.45~s and 8,857 in 593.50~s.
The refinement-time ratios are 18.6 and 5.7 respectively.
These exclude the shared initialization, which takes 63.01 and 179.03~s.
Adding that initialization to both methods gives approximately 4.72-fold and 2.73-fold ratios, rather than the refinement-only ratios.

\begin{table}[tb]
\centering
\caption{RNASim refinement from a shared starting topology. VFT denotes 64-thread VeryFastTree minimum-evolution refinement. P2V denotes minibatch refinement with full-profile acceptance on one RTX A6000. Times exclude the shared initialization.}
\label{tab:rnasim_summary}
\small
\begin{tabular}{@{}rrrrrrr@{}}
\toprule
$N$ & Initial & \multicolumn{2}{c}{Wrong splits} & \multicolumn{2}{c}{Refinement (s)} & Init. (s)\\
\cmidrule(lr){3-4}\cmidrule(lr){5-6}
 & wrong splits & VFT & P2V & VFT & P2V & \\
\midrule
1,000 & 214 & 178 & 176 & 6.88 & 0.06 & 0.58\\
5,000 & 890 & 703 & 696 & 33.38 & 0.68 & 3.40\\
10,000 & 1,667 & 1,313 & 1,320 & 65.98 & 1.55 & 7.82\\
50,000 & 6,729 & 4,980 & 4,950 & 313.53 & 16.84 & 63.01\\
100,000 & 12,375 & 8,857 & 8,802 & 593.50 & 104.45 & 179.03\\
\bottomrule
\end{tabular}
\end{table}

The topological advantage is not uniform.
At $N=10{,}000$, the minibatch result has 1,320 incorrect splits, compared with 1,313 for VeryFastTree.
A full-tree run with eigenvalue-based distances reaches 1,309 in 4.57~s, but uses a different distance estimate.
The full-tree JC69 run returns 1,320 in 3.71~s.
These variants should therefore not be combined into a single claim of consistently better topological accuracy.

At $N=1{,}000$, FastME applied to the GPU-computed JC69 distance matrix returns 177 incorrect splits in 21.68~s with its own BIONJ initialization.
Its shared-start refinement returns 178 in 11.27~s.
The corresponding full-tree GPU result is 176 in 0.01~s, and minibatch refinement returns 176 in 0.06~s.
The experiments differ in initialization and timing scope, so the ratios between their raw elapsed times are not interchangeable.

\section{Discussion}

The component identity provides the mathematical basis for evaluating local skeleton changes against a global BME objective.
The unchanged internal terms form a constant, and the remaining terms factor through balanced component profiles.
This principle is established BME structure.
Its role here is to make explicit which data a minibatch evaluator must preserve and to connect reduced skeleton optimization to Phylo2Vec coordinates and parallel execution.

The distinction between objective evaluation and search is essential.
The distance matrix may itself be estimated, the search may select only a small subset of candidates, and capped profiles may miss useful proposals.
None of these facts contradicts an exact algebraic evaluation of a specified candidate using complete profiles.
Conversely, a correct algebraic identity does not remove floating-point error or establish global convergence.
Changing leaf labels enlarges the collection of coordinate neighborhoods available to the optimizer, but the tested two-orientation schedule alone is not an exhaustive SPR search.

The main computational tradeoff is between reuse and reconstruction.
Dense tables permit rapid sweeps but require quadratic memory.
Column tiling trades memory for repeated arithmetic.
Skeletons reduce the dimension of the optimization problem, but exact construction of their distances can retain a dependence on the full leaf set.
This becomes particularly visible in nonlinear sequence-distance refinement and in candidates involving large condensed components.
Bilinear embedding benchmarks therefore establish a different scaling regime from the sequence experiments.

The empirical results support faster evaluation and competitive refinement on the tested instances.
They do not establish an asymptotic bound on the number of optimization sweeps, statistical superiority across data-generating models, or a universal numerical error bound.
The largest sequence comparisons have no reported replicate uncertainty.
Further comparisons should use matched distance matrices, identical starting trees, explicit timing boundaries, and repeated runs, with the exact implementation revision attached to each result.

\section{Conclusion}

Phylo2Vec coordinates provide a structured interface to quadratic-time BME move evaluation.
Balanced component profiles extend this evaluator to reduced skeletons while preserving the full-tree change for every allowed skeleton update.
The resulting implementation combines parallel gradient evaluation with minibatch proposals, full-profile acceptance, and sequential maintenance of outside information.
Its performance depends on both the skeleton size and the cost of constructing exact component distances.
The mathematical identities establish what is evaluated, while the experiments characterize the speed and solution quality of the implemented search.

\input{phylo2vec_bme_references}

\clearpage
\appendix
\input{phylo2vec_bme_appendix}

\end{document}
